\documentclass[10pt,a4paper]{amsart}
\usepackage[margin=19mm]{geometry}
\usepackage{amsmath,amssymb,mathtools}
\usepackage[scr=rsfs,cal=euler]{mathalfa}
\usepackage{xfrac}
\usepackage[unicode,hidelinks,bookmarks=false]{hyperref}
\newcommand{\ie}{i.e.\ }
\newcommand{\cf}{cf.\ }
\newcommand{\resp}{respectively\ }
\newcommand{\Subsection}{\subsection}
\providecommand{\nobreakdash}{\nobreak\hbox{-}}
\setlength{\emergencystretch}{2em}
\multlinegap0pt
\usepackage{amssymb}
\usepackage{mathtools}
\newtheorem{lemma}{Lemma}
\newtheorem{theorem}{Theorem}
\newcommand{\C}{\mathbb C}
\newcommand{\fall}[2]{(#1)_{#2}}
\DeclareMathOperator{\spn}{span}
\title{Finite-term recurrences in a generalized Bochner--Krall family}
\author{L.M. Anguas, D. Barrios Rolan\'ia, B. Shapiro, M. Tater}
\date{Source: arXiv:2608.21802v1 (CC BY 4.0)}
\begin{document}
\maketitle

\noindent\emph{Source and licence.} Finite-term recurrences in a generalized Bochner--Krall family. Version arXiv:2608.21802v1 (CC BY 4.0), distributed by arXiv under the Creative Commons Attribution 4.0 International licence, http://creativecommons.org/licenses/by/4.0/, as recorded in the arXiv licence metadata for this exact version and shown on the abstract page https://arxiv.org/abs/2608.21802v1. Full licence notice: \texttt{licenses/arxiv-2608.21802-CC-BY-4.0.txt}.\par\smallskip

\noindent\textit{Editorial scope.} Counted unit: the article's theorem thm:jge2 (source0.tex lines 443--446). The article's complete original proof of it is delivered with it (source0.tex lines 448--531, one \texttt{proof} environment of 84 lines). No mathematics is written, repaired or re-derived: the statement and the proof are the article's own lines, in the article's own notation, and no retained line is substituted or re-flowed.  Retained uncounted as the source-local context the proof needs: lines 109--111; lines 182--184; lines 357--372; lines 403--416.  Nothing was omitted from any retained span, so the package carries no omission record.  No retained line contains a citation, so the wrapper carries no bibliography and no bibliography entry.  No retained line contains a figure environment, an image-inclusion command or drawing code, so no image is delivered and the package declares no asset dependency.  Editorial formatting only, in addition: an AMS \texttt{amsart} preamble that restates the article's own declarations and macros used by the retained lines, namely the environments \texttt{lemma}, \texttt{theorem} on separate counters, together with the standard packages those lines need.  The retained environments print in this document as Lemma 1, Theorem 1 in order of appearance, whereas the article prints the same items with the numbering of its own full document.  The complete native source of the article as distributed in the arXiv e-print for this exact version is bundled unchanged as \texttt{sources/208277/source0.tex}.
\begin{equation}\label{eq:general-rec}
\alpha_{n,R}=0\quad\text{whenever } n-R\ge N,
\end{equation}

\begin{equation}\label{eq:Tmon}
Tz^n=\fall{n}{j}z^n+\fall{n}{\ell}z^{n-k},
\end{equation}

\begin{lemma}\label{lem:basis}
Let \(0\le c\le\ell-1\).  Then:
\begin{enumerate}
\item \(M_0=0\), and \eqref{eq:Tmon} reads
\begin{equation}\label{eq:bidiagonal}
Te_0=\Lambda_0e_0,
\qquad
Te_i=\Lambda_ie_i+M_ie_{i-1}\quad(i\ge1);
\end{equation}
in particular \(T\mathcal V_c\subseteq\mathcal V_c\).  Also
\(z^k\mathcal V_c\subseteq\mathcal V_c\), with \(z^k e_i=e_{i+1}\).
\item \(P_{c+ik}\in e_i+\spn\{e_0,\dots,e_{i-1}\}\) for every \(i\ge0\);
consequently \(\{P_{c+ik}\}_{i\ge0}\) is a basis of \(\mathcal V_c\).
\item If moreover \(c\ge m\), then \(M_i>0\) for every \(i\ge1\).
\end{enumerate}
\end{lemma}

\begin{lemma}\label{lem:coordinate}
Let \(0\le c\le\ell-1\) and assume that the numbers \(\Lambda_i\)
\((i\ge0)\) are pairwise distinct; this holds automatically if \(c\ge j\),
and also if \(j=1\).  Then \(\psi^{(c)}_R\circ T=\Lambda_R\psi^{(c)}_R\) and
\begin{equation}\label{eq:coordinate}
\psi^{(c)}_R(e_i)=
\begin{cases}
0,&i<R,\\[2pt]
1,&i=R,\\[2pt]
\displaystyle(-1)^{i-R}\prod_{q=R+1}^{i}
\dfrac{M_q}{\Lambda_q-\Lambda_R},&i>R.
\end{cases}
\end{equation}
\end{lemma}

\begin{theorem}\label{thm:jge2}
Let \(2\le j<\ell\) and \(0\le m<\ell\).  Then the monic eigenpolynomials of
\(T\) do not satisfy any recurrence relation with a fixed number of terms.
\end{theorem}

\begin{proof}
Throughout we work on \(\mathcal V_p\) with
\[
p:=\ell-1\ \ (\ge j),
\]
and we abbreviate \(e_i=z^{p+ik}\), \(\Lambda_i=\fall{p+ik}{j}\),
\(M_i=\fall{p+ik}{\ell}\) and \(\psi_R=\psi^{(p)}_R\).  By
Lemma~\ref{lem:basis}, \(M_0=0\) and \(M_i>0\) for \(i\ge1\) (indeed
\(p+ik\ge\ell-1+k\ge\ell\)), and \(\{P_{p+ik}\}_{i\ge0}\) is a basis of
\(\mathcal V_p\).  Since \(p\ge j\), the sequence \((\Lambda_i)_{i\ge0}\) is
strictly increasing, so Lemma~\ref{lem:coordinate} applies; in particular
\(\Lambda_i-\Lambda_0>0\) for \(i\ge1\).

Suppose, for contradiction, that \eqref{eq:general-rec} holds for some
\(N\ge1\).  Multiplication by \(z\) then lowers the index \(R\) by at most
\(N-1\), so multiplication by \(z^k\) lowers it by at most \(k(N-1)\).
Since \(z^k P_{p+ik}\in\mathcal V_p\) and the basis of \(\mathcal V_p\) is a
subset of the basis \(\{P_n\}\) of \(\C[z]\), the coefficient of \(P_p\) in
\(z^k P_{p+ik}\), which is \(\chi(P_{p+ik})\) for
\[
\chi:=\psi_0\circ(\text{multiplication by }z^k)\in\mathcal V_p^{*},
\]
vanishes whenever
\[
p<(p+ik)-k(N-1),
\]
that is, whenever \(ik>k(N-1)\), and consequently whenever \(i\ge N\).
Both \(\chi\) and
\(\sum_{R=0}^{N-1}\chi(P_{p+Rk})\,\psi_R\) therefore take the same value on
every element of the basis \(\{P_{p+ik}\}_{i\ge0}\), whence
\[
\chi=\sum_{R=0}^{N-1}c_R\psi_R,
\qquad c_R:=\chi(P_{p+Rk}).
\]

Evaluate this identity at \(e_i\) with \(i\ge N\).  Since \(z^k e_i=e_{i+1}\),
Lemma~\ref{lem:coordinate} gives
\[
(-1)^{i+1}\prod_{q=1}^{i+1}\frac{M_q}{\Lambda_q-\Lambda_0}
=
\sum_{R=0}^{N-1}c_R(-1)^{i-R}
\prod_{q=R+1}^{i}\frac{M_q}{\Lambda_q-\Lambda_R}.
\]
Dividing by the nonzero quantity
\((-1)^i\prod_{q=1}^{i}M_q/(\Lambda_q-\Lambda_0)\), we obtain
\begin{equation}\label{eq:contradiction}
-\frac{M_{i+1}}{\Lambda_{i+1}-\Lambda_0}
=
\sum_{R=0}^{N-1}\widehat c_R\,\Pi_R(i)
\qquad (i\ge N),
\end{equation}
where the constants \(\widehat c_R\) and the products \(\Pi_R(i)\) are
\[
\widehat c_R
:=(-1)^R c_R\prod_{q=1}^{R}\frac{\Lambda_q-\Lambda_0}{M_q},
\qquad
\Pi_R(i):=\prod_{q=R+1}^{i}\frac{\Lambda_q-\Lambda_0}{\Lambda_q-\Lambda_R}.
\]
Note that \(\widehat c_R\) does not depend on \(i\).

The left-hand side of \eqref{eq:contradiction} tends to \(-\infty\).  Indeed
\(M_{i+1}=\fall{p+(i+1)k}{\ell}\sim(ki)^{\ell}\) and
\(\Lambda_{i+1}-\Lambda_0\sim(ki)^{j}\) as \(i\to\infty\), so the left-hand
side is asymptotic to \(-(ki)^{\ell-j}\), and \(\ell-j\ge1\) by hypothesis.

The right-hand side of \eqref{eq:contradiction}, on the contrary, stays
bounded.  Fix \(R\).  For \(q>R\) we have \(\Lambda_q>\Lambda_R\ge\Lambda_0\),
so every factor of \(\Pi_R(i)\) satisfies
\[
\frac{\Lambda_q-\Lambda_0}{\Lambda_q-\Lambda_R}
=1+\frac{\Lambda_R-\Lambda_0}{\Lambda_q-\Lambda_R}\ \ge\ 1,
\]
and \(\Pi_R(i)\) is nondecreasing in \(i\).  Moreover
\(\Lambda_q=\fall{p+qk}{j}\) is a polynomial in \(q\) of degree \(j\) with
leading coefficient \(k^j>0\), so there is \(C_R>0\) with
\(\Lambda_q-\Lambda_R\ge C_R q^{j}\) for all large \(q\).  Since \(j\ge2\),
\[
\sum_{q>R}\frac{\Lambda_R-\Lambda_0}{\Lambda_q-\Lambda_R}<\infty,
\]
and \(\log(1+x)\le x\) yields
\(\log\Pi_R(i)\le\text{const}<\infty\).  Hence each \(\Pi_R(i)\) increases to
a finite limit, and the right-hand side of \eqref{eq:contradiction} converges
to a finite number as \(i\to\infty\).  This contradiction proves the theorem.
\end{proof}

\end{document}
